\documentclass[10pt,a4paper]{amsart}
\usepackage{amsmath,amssymb,amsthm,graphicx}
% Retained minimal preamble for the Stacks Project chapter varieties.tex.
% Verbatim excerpt of the chapter preamble: Xy-pic setup (native preamble lines 14-18), the native theorem
% declarations with the mathematical macros (native preamble lines 154-217), and the project own no-op
% definitions of the reference/slogan/history display environments (native wrapper.tex lines 1-3) that the
% counted statement uses for its slogan. Copied mechanically; nothing added, removed or redefined.
% For commutative diagrams we use Xy-pic
\usepackage[all]{xy}

% We use 2cell for 2-commutative diagrams.
\xyoption{2cell}
% Theorem environments.
%
\theoremstyle{plain}
\newtheorem{theorem}[subsection]{Theorem}
\newtheorem{proposition}[subsection]{Proposition}
\newtheorem{lemma}[subsection]{Lemma}

\theoremstyle{definition}
\newtheorem{definition}[subsection]{Definition}
\newtheorem{example}[subsection]{Example}
\newtheorem{exercise}[subsection]{Exercise}
\newtheorem{situation}[subsection]{Situation}

\theoremstyle{remark}
\newtheorem{remark}[subsection]{Remark}
\newtheorem{remarks}[subsection]{Remarks}

\numberwithin{equation}{subsection}

% Macros
%
\def\lim{\mathop{\mathrm{lim}}\nolimits}
\def\colim{\mathop{\mathrm{colim}}\nolimits}
\def\Spec{\mathop{\mathrm{Spec}}}
\def\Hom{\mathop{\mathrm{Hom}}\nolimits}
\def\Ext{\mathop{\mathrm{Ext}}\nolimits}
\def\SheafHom{\mathop{\mathcal{H}\!\mathit{om}}\nolimits}
\def\SheafExt{\mathop{\mathcal{E}\!\mathit{xt}}\nolimits}
\def\Sch{\mathit{Sch}}
\def\Mor{\mathop{\mathrm{Mor}}\nolimits}
\def\Ob{\mathop{\mathrm{Ob}}\nolimits}
\def\Sh{\mathop{\mathit{Sh}}\nolimits}
\def\NL{\mathop{N\!L}\nolimits}
\def\CH{\mathop{\mathrm{CH}}\nolimits}
\def\proetale{{pro\text{-}\acute{e}tale}}
\def\etale{{\acute{e}tale}}
\def\QCoh{\mathit{QCoh}}
\def\Ker{\mathop{\mathrm{Ker}}}
\def\Im{\mathop{\mathrm{Im}}}
\def\Coker{\mathop{\mathrm{Coker}}}
\def\Coim{\mathop{\mathrm{Coim}}}

% Boxtimes
%
\DeclareMathSymbol{\boxtimes}{\mathbin}{AMSa}{"02}

%
% Macros for moduli stacks/spaces
%
\def\QCohstack{\mathcal{QC}\!\mathit{oh}}
\def\Cohstack{\mathcal{C}\!\mathit{oh}}
\def\Spacesstack{\mathcal{S}\!\mathit{paces}}
\def\Quotfunctor{\mathrm{Quot}}
\def\Hilbfunctor{\mathrm{Hilb}}
\def\Curvesstack{\mathcal{C}\!\mathit{urves}}
\def\Polarizedstack{\mathcal{P}\!\mathit{olarized}}
\def\Complexesstack{\mathcal{C}\!\mathit{omplexes}}
% \Pic is the operator that assigns to X its picard group, usage \Pic(X)
% \Picardstack_{X/B} denotes the Picard stack of X over B
% \Picardfunctor_{X/B} denotes the Picard functor of X over B
\def\Pic{\mathop{\mathrm{Pic}}\nolimits}
\def\Picardstack{\mathcal{P}\!\mathit{ic}}
\def\Picardfunctor{\mathrm{Pic}}
\def\Deformationcategory{\mathcal{D}\!\mathit{ef}}
\newenvironment{reference}{}{}
\newenvironment{slogan}{}{}
\newenvironment{history}{}{}
\usepackage{cleveref}
\crefname{originallemma}{original lemma}{original lemmas}
\Crefname{originallemma}{Original lemma}{Original lemmas}
\makeatletter
\newlabel{algebra-lemma-characterize-separable-field-extensions}{{030W}{source}{}{sourcepointer-1}{}}
\newlabel{algebra-lemma-characterize-separable-field-extensions@cref}{{[originallemma][0][]030W}{[1][0][]source}{}{}{}}
\makeatother
\crefname{originallemma}{original lemma}{original lemmas}
\Crefname{originallemma}{Original lemma}{Original lemmas}
\makeatletter
\newlabel{algebra-lemma-finite-dimensional-algebra}{{00J6}{source}{}{sourcepointer-2}{}}
\newlabel{algebra-lemma-finite-dimensional-algebra@cref}{{[originallemma][0][]00J6}{[1][0][]source}{}{}{}}
\makeatother
\crefname{originallemma}{original lemma}{original lemmas}
\Crefname{originallemma}{Original lemma}{Original lemmas}
\makeatletter
\newlabel{algebra-lemma-geometrically-reduced-finite-purely-inseparable-extension}{{030V}{source}{}{sourcepointer-3}{}}
\newlabel{algebra-lemma-geometrically-reduced-finite-purely-inseparable-extension@cref}{{[originallemma][0][]030V}{[1][0][]source}{}{}{}}
\makeatother
\crefname{originallemma}{original lemma}{original lemmas}
\Crefname{originallemma}{Original lemma}{Original lemmas}
\makeatletter
\newlabel{algebra-lemma-minimal-prime-reduced-ring}{{00EU}{source}{}{sourcepointer-4}{}}
\newlabel{algebra-lemma-minimal-prime-reduced-ring@cref}{{[originallemma][0][]00EU}{[1][0][]source}{}{}{}}
\makeatother
\crefname{originalproposition}{original proposition}{original propositions}
\Crefname{originalproposition}{Original proposition}{Original propositions}
\makeatletter
\newlabel{algebra-proposition-dimension-zero-ring}{{00KJ}{source}{}{sourcepointer-5}{}}
\newlabel{algebra-proposition-dimension-zero-ring@cref}{{[originalproposition][0][]00KJ}{[1][0][]source}{}{}{}}
\makeatother
\crefname{originallemma}{original lemma}{original lemmas}
\Crefname{originallemma}{Original lemma}{Original lemmas}
\makeatletter
\newlabel{coherent-lemma-flat-base-change-cohomology}{{02KH}{source}{}{sourcepointer-6}{}}
\newlabel{coherent-lemma-flat-base-change-cohomology@cref}{{[originallemma][0][]02KH}{[1][0][]source}{}{}{}}
\makeatother
\crefname{originallemma}{original lemma}{original lemmas}
\Crefname{originallemma}{Original lemma}{Original lemmas}
\makeatletter
\newlabel{coherent-lemma-proper-over-affine-cohomology-finite}{{02O6}{source}{}{sourcepointer-7}{}}
\newlabel{coherent-lemma-proper-over-affine-cohomology-finite@cref}{{[originallemma][0][]02O6}{[1][0][]source}{}{}{}}
\makeatother
\title{The algebra of global functions of a proper scheme over a field and its connected components}
\author{The Stacks Project Authors}
\date{stacks-project:varieties.tex:local-snapshot-20261011; GNU FDL 1.2 or later}
\begin{document}
\maketitle
\noindent\textit{Editorial note.} All delivered mathematical text is the author's own, copied from the pinned Stacks Project chapter file varieties.tex; no statement or proof step was rewritten, reordered or removed.\par\medskip
\noindent\textit{Editorial note.} The last two parts are obtained by the author's explicit 'of course (5) implies (6)' and 'of course (7) implies (8)' steps; both sentences are reproduced as written.\par\medskip
\noindent\textit{Editorial note.} The counted spans carry 7 navigation stubs and no bibliography citation.\par\medskip
\noindent\textit{Editorial note.} The preamble is the chapter's retained minimal preamble (sources/stacks-varieties/retained-preamble.tex lines 1-77), a verbatim excerpt of the native preamble.\par\medskip
\noindent\textit{Editorial note.} The label lemma-proper-geometrically-reduced-global-sections is defined only in varieties.tex, so no cross-chapter restatement had to be weighed for this card.\par\medskip
\noindent\textit{Allowed standard prerequisites (editorial).}\begin{itemize}
\item Cohomology of a proper scheme over an affine base is finite. Cohomology of Schemes, Lemma coherent-lemma-proper-over-affine-cohomology-finite, invoked at proof lines 1954-1956. Supplies that A is a finite dimensional k-algebra.
\item A finite dimensional algebra over a field is a product of local Artinian algebras. Algebra, Lemma algebra-lemma-finite-dimensional-algebra and Proposition algebra-proposition-dimension-zero-ring, invoked at proof lines 1960-1962. Supplies the product decomposition in part (2).
\item Flat base change for the cohomology of the structure sheaf. Cohomology of Schemes, Lemma coherent-lemma-flat-base-change-cohomology, invoked at proof lines 1982-1983 and 1992-1993. Supplies A tensor k-bar = H\textasciicircum{}0 of the base changed scheme, used in parts (4), (5) and (7).
\item Characterization of separable field extensions and of geometrically reduced purely inseparable extensions. Algebra, Lemmas algebra-lemma-characterize-separable-field-extensions and algebra-lemma-geometrically-reduced-finite-purely-inseparable-extension, invoked at proof lines 1986-1988. Supplies the separability conclusion of part (4).
\end{itemize}\medskip
\ifdefined\hypertarget\hypertarget{sourcepointer-1}{}\fi
\noindent\textit{Source reference.} Original lemma 030W: sources/stacks-chow/tags.txt, source lines 3963--3963. Pointer outside the counted span, invoked once for context; the load-bearing argument of the counted statement is the author's proof reproduced in full in the proof span.\par\medskip
\ifdefined\hypertarget\hypertarget{sourcepointer-2}{}\fi
\noindent\textit{Source reference.} Original lemma 00J6: sources/stacks-chow/tags.txt, source lines 704--704. Pointer outside the counted span, invoked once for context; the load-bearing argument of the counted statement is the author's proof reproduced in full in the proof span.\par\medskip
\ifdefined\hypertarget\hypertarget{sourcepointer-3}{}\fi
\noindent\textit{Source reference.} Original lemma 030V: sources/stacks-chow/tags.txt, source lines 3962--3962. Pointer outside the counted span, invoked once for context; the load-bearing argument of the counted statement is the author's proof reproduced in full in the proof span.\par\medskip
\ifdefined\hypertarget\hypertarget{sourcepointer-4}{}\fi
\noindent\textit{Source reference.} Original lemma 00EU: sources/stacks-chow/tags.txt, source lines 548--548. Pointer outside the counted span, invoked once for context; the load-bearing argument of the counted statement is the author's proof reproduced in full in the proof span.\par\medskip
\ifdefined\hypertarget\hypertarget{sourcepointer-5}{}\fi
\noindent\textit{Source reference.} Original proposition 00KJ: sources/stacks-chow/tags.txt, source lines 753--753. Pointer outside the counted span, invoked once for context; the load-bearing argument of the counted statement is the author's proof reproduced in full in the proof span.\par\medskip
\ifdefined\hypertarget\hypertarget{sourcepointer-6}{}\fi
\noindent\textit{Source reference.} Original lemma 02KH: sources/stacks-chow/tags.txt, source lines 3372--3372. Pointer outside the counted span, invoked once for context; the load-bearing argument of the counted statement is the author's proof reproduced in full in the proof span.\par\medskip
\ifdefined\hypertarget\hypertarget{sourcepointer-7}{}\fi
\noindent\textit{Source reference.} Original lemma 02O6: sources/stacks-chow/tags.txt, source lines 3505--3505. Pointer outside the counted span, invoked once for context; the load-bearing argument of the counted statement is the author's proof reproduced in full in the proof span.\par\medskip
\begin{lemma}
\label{lemma-proper-geometrically-reduced-global-sections}
Let $k$ be a field. Let $X$ be a proper scheme over $k$.
\begin{enumerate}
\item $A = H^0(X, \mathcal{O}_X)$ is a finite dimensional $k$-algebra,
\item $A = \prod_{i = 1, \ldots, n} A_i$ is a product of Artinian
local $k$-algebras, one factor for each connected component of $X$,
\item if $X$ is reduced, then $A = \prod_{i = 1, \ldots, n} k_i$
is a product of fields, each a finite extension of $k$,
\item if $X$ is geometrically reduced, then $k_i$ is finite separable
over $k$,
\item if $X$ is geometrically connected, then $A$ is geometrically
irreducible over $k$,
\item if $X$ is geometrically irreducible, then $A$ is geometrically
irreducible over $k$,
\item if $X$ is geometrically reduced and geometrically connected, then
$A = k$, and
\item if $X$ is geometrically integral, then $A = k$.
\end{enumerate}
\end{lemma}

\begin{proof}
By Cohomology of Schemes, Lemma
\ref{coherent-lemma-proper-over-affine-cohomology-finite}
we see that $A = H^0(X, \mathcal{O}_X)$ is a finite dimensional
$k$-algebra. This proves (1).

\medskip\noindent
Then $A$ is a product of local Artinian $k$-algebras by
Algebra, Lemma \ref{algebra-lemma-finite-dimensional-algebra} and
Proposition \ref{algebra-proposition-dimension-zero-ring}.
If $X = Y \amalg Z$ with $Y$ and $Z$ open in $X$, then we obtain
an idempotent $e \in A$ by taking the section of $\mathcal{O}_X$
which is $1$ on $Y$ and $0$ on $Z$. Conversely, if $e \in A$
is an idempotent, then we get a corresponding decomposition of $X$.
Finally, as $X$ has a Noetherian underlying topological space
its connected components are open. Hence the connected components
of $X$ correspond $1$-to-$1$ with primitive idempotents of $A$.
This proves (2).

\medskip\noindent
If $X$ is reduced, then $A$ is reduced. Hence the local rings $A_i = k_i$
are reduced and therefore fields (for example by
Algebra, Lemma \ref{algebra-lemma-minimal-prime-reduced-ring}).
This proves (3).

\medskip\noindent
If $X$ is geometrically reduced, then
$A \otimes_k \overline{k} =
H^0(X_{\overline{k}}, \mathcal{O}_{X_{\overline{k}}})$
(equality by Cohomology of Schemes, Lemma
\ref{coherent-lemma-flat-base-change-cohomology}) is reduced.
This implies that $k_i \otimes_k \overline{k}$ is a product
of fields and hence $k_i/k$ is separable for example by
Algebra,
Lemmas \ref{algebra-lemma-characterize-separable-field-extensions} and
\ref{algebra-lemma-geometrically-reduced-finite-purely-inseparable-extension}.
This proves (4).

\medskip\noindent
If $X$ is geometrically connected, then $A \otimes_k \overline{k} =
H^0(X_{\overline{k}}, \mathcal{O}_{X_{\overline{k}}})$
is a zero dimensional local ring by part (2) and hence its
spectrum has one point, in particular it is irreducible.
Thus $A$ is geometrically irreducible. This proves (5).
Of course (5) implies (6).

\medskip\noindent
If $X$ is geometrically reduced and geometrically connected, then
$A = k_1$ is a field and the extension $k_1/k$ is finite separable and
geometrically irreducible. However, then $k_1 \otimes_k \overline{k}$
is a product of $[k_1 : k]$ copies of $\overline{k}$ and we conclude
that $k_1 = k$. This proves (7). Of course (7) implies (8).
\end{proof}

\end{document}
